\section*{الفصل \ref{ch:b1:s-block} --- الهيدروجين وعناصر الفئة s}

\begin{solution}{exo:b1:s-block:1}
\ce{KH} و\ce{CaH2}: ملحيان؛ \ce{SiH4} و\ce{H2S}: تساهميان؛
\ce{PdH_{0.6}}: فلزي. \ce{CaH2 + 2H2O -> Ca(OH)2 + 2H2}.
\end{solution}

\begin{solution}{exo:b1:s-block:2}
\ce{2Li + 2H2O -> 2LiOH + H2}، والتفاعل نفسه مع Na وK؛
\ce{Mg + 2HCl -> MgCl2 + H2}.
\end{solution}

\begin{solution}{exo:b1:s-block:3}
قاعدية: \ce{Na2O}، \ce{MgO} (\ce{MgO + 2HCl -> MgCl2 + H2O}). حمضية:
\ce{SiO2}، \ce{SO3}، \ce{CO2} (\ce{SO3 + H2O -> H2SO4}). مذبذبة:
\ce{Al2O3}، \ce{ZnO} (\ce{ZnO + 2OH- + H2O -> Zn(OH)4^2-}).
\end{solution}

\begin{solution}{exo:b1:s-block:4}
عند pH يساوي 14 تكون مزدوجة الماء \ce{2H2O + 2e- -> H2 + 2OH-}، \babelsublr{$E = -0.83$~V}؛
$\log K = 2(-0.83 + 2.71)/0.059 = 64$.
\end{solution}

\begin{solution}{exo:b1:s-block:5}
في اللهب، ترفع الاصطدامات إلكترونًا إلى مستوى مثار؛ وعند عودته
إلى \omterm{def:b1:quantum-numbers:energy-level}{الحالة الأساسية}، تصدر الذرة فوتونًا طاقته تساوي الفجوة،
وهي مميزة للعنصر. $E = hc/\lambda = 6.626 \times 10^{-34} \times
3.00 \times 10^8/(589 \times 10^{-9}) = \qty{3.375e-19}{J} = \qty{2.11}{eV}$.
% ledger: const:h, const:c, const:eV
\end{solution}

\begin{solution}{exo:b1:s-block:6}
$10^6/106.0 = \qty{9.43e3}{mol}$ من رماد الصودا: \qty{9.43e3}{mol} من
الحجر الجيري، أي \qty{0.944}{t}، و$2 \times 9.43 \times 10^3/0.90 =
\qty{2.10e4}{mol}$ من الملح، أي \qty{1.23}{t}.
\end{solution}

\begin{solution}{exo:b1:s-block:7}
$[\ce{Mg^2+}] = 1.3/24.3 = \qty{0.053}{mol/L}$؛ $\mathrm{pH} = 14.00 -
\frac12(11.25 + \log 0.053) = 9.0$. و\qty{53.5}{mol} من المغنيزيوم في
\qty{1.0}{m^3} تحتاج إلى \qty{53.5}{mol} من \ce{Ca(OH)2}: أي \qty{4.0}{kg}.
\end{solution}

\begin{solution}{exo:b1:s-block:8}
ثابت التوازن (الديناميكا الحرارية) هائل في الحالتين؛ أما السرعة
(الحركية) فلا يحسمها $E^\circ$. فالليثيوم ينصهر عند درجة حرارة
أعلى، فيبقى صلبًا ويغطيه طبقة ناتج أقل ذوبانًا،
ويفقد إلكترونه بسهولة أقل من البوتاسيوم (طاقة تأين
أعلى): فيتفاعل بانتظام، والبوتاسيوم بعنف.
\end{solution}

\begin{solution}{exo:b1:s-block:9}
$10^6/100.1 = \qty{9.99e3}{mol}$: CaO $9.99 \times 10^3 \times 56.1 =
\qty{0.560}{t}$؛ و\ce{CO2} \qty{9.99e3}{mol}، $V = nRT/p = 9.99 \times 10^3
\times 8.314 \times 298.15/10^5 = \qty{248}{m^3}$.
\end{solution}

\begin{solution}{exo:b1:s-block:10}
$2.9 \times 10^8/6.9 = \qty{4.2e7}{kmol}$ من Li (والكتلة بالكيلوغرام)، $\times 73.8/2 =
\qty{1.55e9}{kg}$: أي نحو 1.6 مليون طن من كربونات الليثيوم.
$2.9 \times 10^8/8 = \num{3.6e7}$ بطارية.
\end{solution}

\begin{solution}{exo:b1:s-block:11}
1.31 و0.84 و0.71 بالمئة عند 20 و80 و\qty{100}{\celsius}: \omterm{def:b1:precipitation:solubility}{فالذوبانية}
تتناقص. وعلى الساخن، يبقى من الكربونات في المحلول قدر أقل: فيُسترجع قدر أكبر.
\end{solution}

\begin{solution}{exo:b1:s-block:12}
\ce{NaHCO3} هو الأقل ذوبانًا بين الأملاح التي يمكن أن تكوّنها الأيونات،
ويتشبّع المحلول به أولًا. والأمونيا تجعل المحلول قاعديًا
بما يكفي لتحويل \ce{CO2} إلى \ce{HCO3-} مع إمكان استرجاعها (على هيئة
\ce{NH4Cl}، الذي يجدده الجير)؛ أما هيدروكسيد الصوديوم فسيُستهلك
ويكلّف أكثر من الناتج.
\end{solution}

\begin{solution}{pb:b1:s-block:1}
\textbf{1.} Li \qty{6.0}{g/L}، Mg \qty{2.4}{g/L}.
\textbf{2.} أربعة أمتار مكعبة من المحلول الملحي تعطي مترًا واحدًا: فيتبخر ثلاثة أمتار مكعبة من
الماء.
\textbf{3.} ضوء الشمس والمناخ الجاف يوفّران الطاقة مجانًا.
\textbf{4.} كلوريد الصوديوم، وهو أوفر الأملاح إلى حد بعيد، يتشبّع
أولًا.
\textbf{5.} Li $6000/6.9 = \qty{870}{mol}$؛ Mg $2400/24.3 = \qty{98.8}{mol}$.
\textbf{6.} كان المغنيزيوم سيترسب مع الكربونات
ويلوّث الناتج.
\textbf{7.} \ce{Mg^2+ + Ca(OH)2 -> Mg(OH)2 + Ca^2+}.
\textbf{8.} $98.8 \times 74.1 = \qty{7.32}{kg}$.
\textbf{9.} يبقى $[\ce{Mg^2+}] = \num{9.88e-5}$ mol/L: $\mathrm{pH} = 14.00 -
\frac12(11.25 - 4.01) = 10.38$.
\textbf{10.} هيدروكسيد الليثيوم ذوّاب: فلا يتكوّن صلب عند قيمة pH هذه.
\textbf{11.} \ce{Ca^2+ + CO3^2- -> CaCO3}.
\textbf{12.} $98.8 \times 106.0 = \qty{10.5}{kg}$.
\textbf{13.} \ce{2Li+ + CO3^2- -> Li2CO3}.
\textbf{14.} $435 \times 106.0 = \qty{46.1}{kg}$.
\textbf{15.} $435 \times 73.8 = \qty{32.1}{kg}$.
\textbf{16.} كربونات الليثيوم أقل ذوبانًا على الساخن.
\textbf{17.} $0.0084 \times 1000 = \qty{8.4}{kg}$.
\textbf{18.} بكميات صغيرة من الماء الساخن، الذي تكون فيه أقل
ذوبانًا.
\textbf{19.} $32.1 - 8.4 = \qty{23.7}{kg}$.
\textbf{20.} $23.7/32.1 = \qty{74}{\percent}$.
\textbf{21.} ما زال يحتوي ليثيومًا: فيُعاد إلى الأحواض.
\textbf{22.} $2.9 \times 10^8/6.0 = \num{4.8e7}$ متر مكعب.
\textbf{23.} \textbf{\qty{24}{kg}} (23.7) من كربونات الليثيوم لكل متر
مكعب.
\end{solution}
